\section{Introduction}\label{sec:intro}

Forecasting daily stock index levels is a problem where small reported gains
attract a great deal of attention, because a model that is even slightly more
accurate than its competitors has obvious value to anyone managing risk.
Machine learning for stock market prediction is by now a large literature, with
well over a hundred journal articles reviewed in a single survey~\cite{kumbure2022},
and over the past decade one design has come to dominate the applied part of
it: split the price series into simpler subseries with a signal decomposition,
forecast each subseries with its own neural network, and add the forecasts back
together. Empirical mode decomposition (EMD) and its interpolation and noise
assisted variants~\cite{akimaemd2023,gruceemdan2023,ceemdanstock2025} and
variational mode decomposition (VMD)~\cite{jiang2025,wangwang2026} are the usual
front ends, and recurrent and convolutional networks, often with attention, are
the usual back ends~\cite{gong2024,zhao2025,ning2025}. Reported accuracies are
often high, with coefficients of determination close to
one~\cite{ning2025,jiang2025}.

Numbers of that kind deserve a second look for a specific reason. A signal
decomposition is not a causal filter: the value of a component at time $t$
depends on the whole stretch of data that was decomposed, including samples
after $t$. When a study decomposes the complete series first and only then
separates training from test data, every input the model sees at test time has
been computed with knowledge of the prices it is asked to predict. Leakage of
this general kind has produced overoptimistic conclusions in many fields that
adopted machine learning~\cite{kapoor2023}. For decomposition based
forecasters it is well documented in wind power forecasting, where it is known
as the boundary issue~\cite{chen2022}; it has been shown to create an illusion
of accuracy in water quality prediction~\cite{emdleak2024}; and several recent
studies design decomposition schemes that keep future samples out of the model
inputs~\cite{wang2024,yao2025,pang2026}. Look ahead bias is also a recognised
concern in financial backtesting, most recently for signals extracted by large
language models~\cite{glasserman2024}, and a few stock forecasting studies have
replaced full series decomposition with sliding window or rolling decomposition
for exactly this reason~\cite{liu2022,liu2024,wen2024,zhang2025}. What we could not
find for equity indices is a measurement of how much of the accuracy of an
existing model the leak accounts for, or a test of whether such a model still
beats persistence once the leak is removed. That question matters because,
under the random walk and efficient market hypotheses, future prices cannot be
systematically predicted~\cite{pernagallo2025}, even though machine learning
classifiers have been reported to beat random choice on longer horizon relative
returns~\cite{finrel2024}.

This paper measures how much of the apparent skill of a decomposition ensemble
stock forecaster survives when the decomposition is made strictly causal. We
work with a published redecomposition model that combines improved complete
ensemble EMD with adaptive noise (ICEEMDAN), a second VMD stage tuned by
particle swarm optimisation (PSO), and a bidirectional long short term memory
(BiLSTM) network with self attention feeding a temporal convolutional network
(TCN)~\cite{gong2024}, together with a variant that replaces the convolutional
head with a Tiny Recursive Model (TRM)~\cite{jolicoeur2025}. On the daily
highest and lowest prices of the S\&P~500 and the Shanghai Stock Exchange
Composite Index (SSEC) we compare two protocols that differ in one respect
only: whether the series is decomposed once in full or redecomposed from the
available history at every forecast origin. The contributions are as follows.
\begin{enumerate}
\item We show, before any network is trained, that full series decomposition
places information about the next day's price change into the model inputs,
and that walk forward decomposition removes it.
\item We quantify the accuracy attributable to decomposing the full series,
as the ratio of out of sample error under walk forward decomposition to the
error under full series decomposition, measured on identical test dates for
every index, price series and network.
\item We evaluate both forecasters against persistence, drift, moving average
and exponential smoothing benchmarks with the modified Diebold and Mariano
(MDM) test~\cite{hewamalage2022,mdmindex2023}, so that any claimed advantage is
judged against forecasts that require no model at all.
\item In reproducing the pipeline we identified fourteen issues. Twelve were
corrected, several of which would on their own have changed the central
comparison, and two are modelling choices of the source model that we keep
and report. Section~\ref{sec:audit} documents each one, so that the
reproduction can be checked rather than trusted.
\item We give both network heads identical training budgets, schedules and
postprocessing, which the source comparison did not, and report whether the
ranking of the two heads depends on the protocol.
\end{enumerate}
Section~\ref{sec:related} reviews the related work, Section~\ref{sec:data}
describes the data and the two protocols, Section~\ref{sec:setup} the
networks, training, measures and the audit of the reproduced pipeline,
Section~\ref{sec:results} the results, Section~\ref{sec:discussion} their
interpretation and limits, and Section~\ref{sec:conclusion} concludes.

\section{Related work and its limitations}\label{sec:related}

This section covers three topics in turn: the decomposition ensemble design
and its recent extensions, the evidence that a decomposition can carry future
information into a forecaster, and the benchmarks and network heads that this
study uses.

\subsection{Decomposition ensemble forecasting}
EMD represents a nonstationary series as a sum of intrinsic mode functions
extracted by an iterative, spline based sifting procedure whose behaviour at
the boundaries of the data needs particular care~\cite{emdtutorial2023}. Plain
EMD is prone to mode mixing, and ICEEMDAN is one of the variants adopted to
counter it~\cite{iceemdanbridge2022}; such noise assisted methods decompose
noise added copies of the signal, together with the added noise itself, using
classical EMD~\cite{ceemdanpaa2022}. VMD extracts a preset number of modes and
cannot choose its input parameters by itself~\cite{vmdmssr2022}, so they are
frequently tuned by an optimiser such as PSO~\cite{psovmdtcn2023}.
Decomposition based hybrids have become very popular~\cite{chen2022}, and
recent work continues to extend them: EMD with long short term memory (LSTM)
networks for the KSE-100 index~\cite{akimaemd2023}, complete ensemble EMD with
adaptive noise (CEEMDAN) with wavelet denoising and gated recurrent units for
the S\&P~500 and CSI~300~\cite{gruceemdan2023}, CEEMDAN with LSTM and
backpropagation networks for the DAX~\cite{ceemdanstock2025}, a second VMD
stage applied to the high frequency modes of a CEEMDAN decomposition for wind
power~\cite{ning2025}, VMD followed by convolutional, bidirectional recurrent
and attention layers for natural gas futures prices~\cite{jiang2025}, adaptive
VMD with feature selection and a BiLSTM for stock prices~\cite{wangwang2026},
PSO tuned VMD with an attention TCN for electricity load~\cite{psovmdtcn2023},
wavelet denoising combined with an attention BiLSTM, an autoregressive
integrated moving average (ARIMA) model and gradient boosting for the CSI~300
index~\cite{zhao2025}, metaheuristic tuning of decomposition aided attention
recurrent networks for solar and wind generation~\cite{damasevicius2024}, and
decomposition blocks placed inside general forecasting
architectures~\cite{kreuzer2025}.

\subsection{Future information in decomposition}
Chen \emph{et al.} reviewed decomposition based wind power models and noted
that existing studies usually decompose data that would not be available at
forecast time, so the reported accuracy is higher than what can be achieved in
practice~\cite{chen2022}. Yang \emph{et al.} found the same effect in EMD
based water quality prediction, where decomposing before the split let test
information into training and produced an illusion of accuracy, and they
proposed sliding window and repeated decomposition schemes to remove
it~\cite{emdleak2024}. Other remedies include cascade decomposition designed to
avoid information leakage~\cite{wang2024}, rolling decomposition in which the
point to be predicted never enters the decomposition~\cite{yao2025}, and
explicit guidance on preprocessing that prevents future data from leaking into
wind speed forecasts~\cite{pang2026}. In stock forecasting, Liu \emph{et al.}
decompose sliding windows rather than the full series to remove the
leak~\cite{liu2022,liu2024}, Wen \emph{et al.} design a wavelet network that avoids
it~\cite{wen2024}, and Zhang \emph{et al.} apply rolling CEEMDAN to realised
volatility~\cite{zhang2025}. The walk forward protocol used here follows the
same principle and matches recommended practice for partitioning time series
in forecast evaluation~\cite{hewamalage2022}. These studies propose designs that
avoid the leak and report the accuracy of the new design. We found none that
measures, on equity indices, how much of the accuracy of an existing published
model the leak accounts for under otherwise identical training, that tests the
causal version against persistence, or that shows the leak in the inputs before
any network is trained; those are the gaps this paper addresses.

\subsection{Benchmarks and network heads}
Under the random walk and efficient market hypotheses future index levels
cannot be systematically predicted~\cite{pernagallo2025}, so persistence is a
demanding benchmark, and simple models deserve a place in any comparison: a one
layer linear model outperformed sophisticated Transformer based forecasters on
nine real datasets~\cite{zeng2023}, and a large comparison on the M3 competition
series found that deep learning adds value mainly in combination and at a
considerable computational cost~\cite{makridakis2022}. Machine learning
studies often adopt evaluation practices that make uncompetitive methods look
competitive, and we follow published guidance on data partitioning, error
measures and statistical testing~\cite{hewamalage2022}; a systematic survey of
deep learning in the stock market likewise calls for more reproducible,
explainable and accountable models~\cite{olorunnimbe2022}. Forecasts are
compared with the MDM test, which has been used to compare daily index price
forecasts~\cite{mdmindex2023}. The encoder combines a BiLSTM with attention, a
combination covered in recent reviews of recurrent
architectures~\cite{rnnreview2024} and the one used by the source
model~\cite{gong2024}. One head is a TCN, which recent decomposition based
forecasters also pair with attention~\cite{psovmdtcn2023,emdleak2024}; the
other adapts the recursive latent refinement of the TRM~\cite{jolicoeur2025}.

\section{Data and decomposition protocols}\label{sec:data}

The experiment is defined by the price series, the two stage decomposition
applied to them, and the two protocols under which that decomposition is run.
Five subsections describe them: the series, the decomposition, protocol A,
protocol B, and how results under the two protocols are matched.

\subsection{Price series}
The data are the daily highest and lowest prices of the S\&P~500 (ticker
\texttt{\textasciicircum GSPC}) and of the SSEC (\texttt{000001.SS}) over a
ten year window, downloaded once as unadjusted values and cached, so that
every experiment uses the same array. Rows with a missing highest or lowest
price were removed jointly, keeping the two series of an index aligned by
trading day. The S\&P~500 series contain $2513$ observations and the SSEC
series $2425$. Each series is forecast one step ahead from a window of the
previous $30$ values. Of the windowed samples, the first $80\%$ form the
training set and the rest the test set, giving $497$ test days for the
S\&P~500 and $479$ for the SSEC.

\subsection{Two stage decomposition}
The first stage applies ICEEMDAN~\cite{colominas2014} to the price
series, producing intrinsic mode functions ordered from highest to lowest
frequency and a final residual. Sifting stops when the normalised squared
difference between successive iterates falls below $0.2$, with at most $100$
sifting iterations and $10$ modes. The second stage decomposes the highest
frequency mode again with VMD~\cite{dragomiretskiy2014}, with the number of modes $K$
and the bandwidth penalty $\alpha$ chosen by PSO~\cite{psosurvey2022} over
$K\in[3,10]$ and $\alpha\in[100,5000]$ (twenty particles, thirty iterations,
inertia $0.7$, both acceleration coefficients $1.5$). The $K$ VMD modes, the
remaining ICEEMDAN modes and the residual form the $C$ subseries that are
forecast, with $C$ equal to $16$ or $17$ here.

Two properties of this pipeline, inherited from the source model and kept
unchanged so that the reproduction stays faithful, matter for the
interpretation of the results. First, the swarm minimises the reconstruction
error of the decomposed mode, and that error falls monotonically as $K$ grows
and $\alpha$ shrinks. The search therefore returned the bounds $K=10$ and
$\alpha=100$ on all eight decompositions, so the optimisation step does not in
practice select anything. Second, the noise added in each ICEEMDAN realisation
has a fixed standard deviation of $0.2$ price units, between
$2\times10^{-4}$ and $8\times10^{-4}$ of the standard deviation of the
training series. This amplitude is small, although it still alters the
sifting: two runs that differ only in their random seed disagree by up to
$1.5\%$ of the price range in individual components. Our results should
therefore be read as describing ICEEMDAN at this noise level.

\subsection{Protocol A: full series decomposition}
Let $x_1,\dots,x_N$ be a price series, $\mathcal{D}$ the two stage
decomposition, and $c_k(t)$ the value of component $k$ at time $t$. A
forecast of $x_t$ is formed from component forecasts,
\begin{equation}
\hat{x}_t=\sum_{k=1}^{C}\hat{c}_k(t),\qquad
\hat{c}_k(t)=f_k\bigl(c_k(t-30),\dots,c_k(t-1)\bigr),
\label{eq:ensemble}
\end{equation}
where each $f_k$ is a separately trained network. The protocols differ only in
how the component values entering $f_k$ are computed. Under protocol A the
whole series is decomposed once, $c_k(\cdot)=\mathcal{D}_k(x_{1:N})$, and the
resulting components are then divided into training and test windows. Both the
swarm search and every component value, including those inside test windows,
depend on all $N$ observations. This is the protocol of the source model and of
much of the literature reviewed above.

\subsection{Protocol B: walk forward decomposition}
Under protocol B the training portion alone is decomposed, and the swarm
search runs on it alone; the resulting $K$, $\alpha$ and $C$ are then held
fixed. For every test origin $t$ the available history is decomposed afresh
and the last $30$ rows of that decomposition form the input,
\begin{equation}
c_k^{(t)}(s)=\mathcal{D}_k\bigl(x_{a_t:t-1}\bigr)(s),\qquad s\le t-1,
\label{eq:walkforward}
\end{equation}
so no value at or after $t$ is ever used to build an input. VMD needs an input
of even length, and $a_t\in\{1,2\}$ is chosen to provide one: when the history
has odd length its oldest observation is left out, so the newest observation
$x_{t-1}$ is always kept. Decomposing histories of different lengths can yield
a different number of intrinsic mode functions; surplus low frequency columns
are added to the residual and a shortfall is filled with a zero column
immediately before it, which leaves the sum of the components unchanged. Each
series needs one decomposition per test day plus one, $498$ for the S\&P~500
and $480$ for the SSEC.

\subsection{Matching the two protocols}
Two practical differences between the protocols were measured rather than
assumed away. Protocol A was run with $50$ noise realisations, as in the
source model, and protocol B with $20$ to keep the $1956$ walk forward
decompositions tractable. Decomposing the training portion of the highest
price series both ways, the largest difference in any component was $1.3\%$
of the price range, comparable to the $1.5\%$ separating two realisation runs
of $50$ that differ only in their seed, so the smaller ensemble adds about as
much variation as a change of seed and no more. In addition, the full series
decomposition trims the series to an even length at its end, so the two
protocols place their split one day apart. All comparisons between protocols
are therefore computed on the dates both protocols test, $496$ for the
S\&P~500 and $478$ for the SSEC.

\section{Experimental setup}\label{sec:setup}

Five subsections describe the setup: the two networks, their training, the
evaluation measures, the issues identified in the reproduced pipeline and how
each was resolved, and what a complete run of the study consists of.

\subsection{Networks}
Both forecasters share an encoder: a two layer BiLSTM with $64$ units per
direction and dropout of $0.1$ between layers~\cite{rnnreview2024}, followed by
four head self attention with a residual connection and layer normalisation,
as in the source model~\cite{gong2024}. The \emph{TCN head} applies a causal
dilated convolution (kernel $3$, dilation $2$) with a rectified linear unit,
dropout of $0.1$ and a residual linear path, and reads the forecast from the
last position. The \emph{TRM head} projects the last encoder state and refines
a latent state $z$ and an answer state $y$ recursively, applying a two layer
latent network six times per cycle and an answer network once, for three
cycles of which only the last is differentiated~\cite{jolicoeur2025}.

\subsection{Training}
One network is trained per component, so each protocol, index, price series
and head involves $16$ or $17$ networks and the study trains $262$ in total.
Each input window is expressed relative to its own last value, and the network
$g_k$ forecasts the change from that value to the next one,
\begin{equation}
\tilde{u}_j=\frac{c_k(t-31+j)-c_k(t-1)}{\sigma_{\mathrm{in}}},\quad j=1,\dots,30,
\qquad
\hat{c}_k(t)=c_k(t-1)+\sigma_{\mathrm{out}}\,g_k(\tilde{u}),
\label{eq:scaling}
\end{equation}
where $\sigma_{\mathrm{in}}$ and $\sigma_{\mathrm{out}}$ are the standard
deviations of the input deviations and of the one step changes over the
training windows. Component forecasts from~\eqref{eq:scaling} are then summed
as in~\eqref{eq:ensemble}. Section~\ref{sec:audit} explains why this replaces
the min max scaling of the source model. The two heads receive identical
treatment: the Adam optimiser with learning rate $5\times10^{-4}$ and weight
decay $10^{-5}$, minibatches of $64$, gradient clipping at norm $1$, up to
$150$ epochs with early stopping after $20$ epochs without improvement in the
training loss, and the learning rate halved after $8$ such epochs. Because
stopping is decided on the training loss, no data outside the training windows
influence any network. The weights used for prediction are an exponential
moving average with decay $\beta=0.999$, corrected for its initialisation,
\begin{equation}
\bar{\theta}_T=\frac{(1-\beta)\sum_{t=1}^{T}\beta^{\,T-t}\theta_t}{1-\beta^{T}},
\label{eq:ema}
\end{equation}
so that the average is taken over the weights actually visited and does not
depend on how long a component trained. Seeds are fixed per component. All
networks were trained on a single NVIDIA GeForce GTX~1650 graphics card.

\subsection{Evaluation measures}
Accuracy is reported as root mean squared error (RMSE), mean absolute error
(MAE), mean absolute percentage error (MAPE) and the coefficient of
determination $R^2$, following guidance on choosing error measures for
forecast evaluation~\cite{hewamalage2022}. Five benchmarks need no learning:
persistence, $\hat{x}_t=x_{t-1}$; a random walk with drift estimated on the
training data; moving averages of the last $5$ and $20$ values; and
exponential smoothing with span $10$. Two forecasts $a$ and $b$ are compared
with the MDM test~\cite{mdmindex2023}. With squared error loss differential
$d_t=e_{a,t}^2-e_{b,t}^2$ over $T$ test days and horizon $h=1$, the statistic is
\begin{equation}
\mathrm{MDM}=\sqrt{\frac{T+1-2h+h(h-1)/T}{T}}\;
\frac{\bar{d}}{\sqrt{\hat{\gamma}_0/T}},
\label{eq:mdm}
\end{equation}
where $\bar{d}$ is the mean and $\hat{\gamma}_0$ the variance of $d_t$, and
the statistic is referred to a Student $t$ distribution with $T-1$ degrees of
freedom. A positive value means that $a$ has the larger loss. The effect of the
protocol is summarised by the leakage inflation ratio
\begin{equation}
\rho=\frac{\mathrm{RMSE}_{\mathrm{B}}}{\mathrm{RMSE}_{\mathrm{A}}},
\label{eq:rho}
\end{equation}
computed on common test dates for each index, price series and head. A value
of $\rho=1$ means that decomposing the full series gave no advantage; a value
above one measures how much of the full series accuracy depends on future
data.

\subsection{Issues identified in the reproduced pipeline}\label{sec:audit}
The reproduction started from the source implementation and changed it only
where a change was needed for the comparison to be valid.
Tables~\ref{tab:audit_decomp} and~\ref{tab:audit_train} list every issue we
identified, its consequence, and how it was resolved, the first for the
decomposition and data handling stages and the second for training, so that a
reader can judge which of the differences between our numbers and the source's
are due to which change. Where the original behaviour was a modelling choice
rather than a fault, it was kept and reported.

\begin{table}[ht]
\caption{Issues identified in the decomposition and data handling stages of the source pipeline, and their resolution.}
\label{tab:audit_decomp}
\small
\begin{tabular}{@{}L{0.30\textwidth}L{0.39\textwidth}L{0.23\textwidth}@{}}
\toprule
Issue & Consequence & Resolution \\
\midrule
ICEEMDAN returned its components in reverse order &
VMD redecomposed the lowest frequency residual instead of the highest frequency mode &
Order corrected \\
Realisations with fewer modes were padded after the residual &
Averaging mixed one realisation's residual with another's mode &
Padding placed before the residual \\
Swarm leader compared against its own overwritten fitness &
The global best could not improve once it led &
Global best tracked separately \\
Walk forward component target copied from the last input value &
Component level scores trivially perfect; price level scores unaffected &
Target taken from the decomposition that includes day $t$ \\
Walk forward test days dropped when the mode count changed &
Up to $96.5\%$ of test days discarded (SSEC highest: $17$ of $479$ kept) &
Mode count reconciled; all test days kept \\
Walk forward histories of odd length trimmed at their newest end &
On alternate test days the input window ended at $t-2$, turning half the walk forward forecasts into two step forecasts &
Oldest observation left out instead, Eq.~\eqref{eq:walkforward}; every window ends at $t-1$ \\
Highest and lowest prices cleaned of missing rows separately &
The two series could refer to different trading days &
Missing rows removed jointly \\
Protocols split one day apart &
Row $i$ of each protocol refers to a different date &
Comparisons on common dates \\
Swarm objective decreases monotonically in $K$ and $1/\alpha$ &
Search returns the bounds $K=10$, $\alpha=100$ on every run &
Kept as published; reported \\
ICEEMDAN noise fixed at $0.2$ price units &
Noise is $2$ to $8\times10^{-4}$ of the series standard deviation &
Kept as published; reported \\
\botrule
\end{tabular}
\end{table}

\begin{table}[ht]
\caption{Issues identified in the training stage of the source pipeline, and their resolution.}
\label{tab:audit_train}
\small
\begin{tabular}{@{}L{0.30\textwidth}L{0.39\textwidth}L{0.23\textwidth}@{}}
\toprule
Issue & Consequence & Resolution \\
\midrule
Moving average of the weights initialised at the random initialisation &
After $930$ updates, $39\%$ of the averaged weights were still the random initialisation &
Bias correction, Eq.~\eqref{eq:ema} \\
Heads trained differently (TCN: $50$ epochs at learning rate $10^{-3}$, no early stopping, clipping or averaging; TRM: $150$ epochs at $5\times10^{-4}$ with all three, and a bias correction applied to it alone) &
Head comparison confounded with training treatment &
Identical training for both heads \\
Components min max scaled on their training rows &
The S\&P~500 trend component lay entirely above its training maximum throughout the test period; full series S\&P~500 high forecasts reached RMSE $198$ (TCN) and $915$ (TRM) against $47$ for persistence &
Windows taken relative to their last value, Eq.~\eqref{eq:scaling} \\
TRM head under min max scaling &
Constant forecasts for $11$ of the $16$ S\&P~500 high components &
Resolved by the same change \\
\botrule
\end{tabular}
\end{table}

\subsection{Reproducibility}
A single run of the study consists of five steps, and the counts below are
read from the files each step saves.
\begin{enumerate}
\item Download the unadjusted daily highest and lowest prices of both indices
once, remove rows with a missing value jointly, and cache the four series
($2513$ observations for each S\&P~500 series and $2425$ for each SSEC series).
\item For protocol A, decompose each full series once with $50$ noise
realisations, search $K$ and $\alpha$ with PSO, and window the components,
giving $497$ test days per S\&P~500 series and $479$ per SSEC series.
\item For protocol B, decompose the training portion with $20$ realisations,
fix $K$, $\alpha$ and $C$, and decompose every history from the split onward:
$498$ decompositions per S\&P~500 series and $480$ per SSEC series, $1956$ in
all.
\item Train one network per component, protocol, series and head, each with a
fixed seed: $16$ or $17$ components, two protocols, four series and two heads,
$262$ networks in all, producing $16$ forecast files.
\item Score every forecast file against the price on the dates shared by both
protocols ($496$ per S\&P~500 series and $478$ per SSEC series), together with
the five benchmarks and the MDM test.
\end{enumerate}
Five secondary experiments support the main comparison: the input diagnostic,
which asks whether the inputs alone predict the next day's change before any
network is trained (Table~\ref{tab:diagnostic}, Fig.~\ref{fig:mechanism}); the
realisation count check, which asks whether $20$ realisations in protocol B
change the components relative to $50$ (Section~\ref{sec:data}); the boundary
amplitude comparison, which asks how variable freshly decomposed components are
at the newest end of the history (Section~\ref{sec:results}); the input
persistence check, which asks what a network would score by copying its own
inputs (Section~\ref{sec:discussion}); and the archived run under the source
model's min max scaling, whose outcome is recorded in Table~\ref{tab:audit_train}.
Every table and figure in this paper is generated by script from the saved
forecast files.
